\appendix
\section{Details for the proof of Theorem~\ref{thm:alg,interp}}\label{app:interp,details}
%
What follows is a rather precise description of results developed by Barvinok in~\cite{Barvinok16,Barvinok18}, lightly specialised to our application and notation.
%
\begin{definition}\label{def:taylor}
Let $f:\CC\to\CC$ be any function, and $m\ge 0$. Then we define the \emph{degree-$m$ Taylor polynomial of $f$ about zero}, $T_m(f)$, as the polynomial in $z$ given by
\[
T_m(f)(z) := f(0) + \sum_{k=1}^m\frac{f^{(k)}(0)}{k!}z^k\,.
\]
\end{definition}
%
\begin{lemma}[{see~\cite[Lemma 2.2.1]{Barvinok16} or \cite[Lemma 2.1]{Barvinok18}}]\label{lem:tayloraccuracy}
Let $g:\hCC\to\hCC$ be a polynomial of degree at most $N$, and for $\beta>1$ suppose that $g(z)\ne 0$ for $|z|<\beta$.
%
Then given a choice of branch for $f(z)=\log g(z)$ where $|z|<\beta$, we have 
\[
|f(1)-T_m(f)(1)| \le \frac{N}{(m+1)\beta^m(\beta-1)}\,.
\]
\end{lemma}
%
\begin{corollary}[{cf.~\cite[Corollary 2.2]{Barvinok18}}]\label{cor:tayloreps}
For any $c>0$ there exists $c'>0$ such that the following holds.
Suppose that the conditions of Lemma~\ref{lem:tayloraccuracy} hold, and in addition that $\beta\le 1+c$. 
%
Then for any $0<\alpha<N/e$, and for any 
\[
m \ge \frac{c'}{\beta-1}\log\left(\frac{N}{\alpha}\right)\,,
\]
we have $|f(1)-T_m(f)(1)|\le \alpha$.
\end{corollary}
%
The only differences from~\cite[Corollary 2.2]{Barvinok18} are the relaxation of the assumption $\alpha<1$ to $\alpha\le N/e$, the additional assumption that $\beta$ is close to $1$, and a more precise analysis of $m$. 
In fact one can take $c'=c/\log(1+c)$.
%
\begin{proof}
By Lemma~\ref{lem:tayloraccuracy}, we are done if 
\[
\frac{N}{(m+1)\beta^m(\beta-1)} \le \alpha
\]
for $m$ as in the statement of the corollary. This holds if and only if
\[
(m+1)\beta^{m+1} \ge \frac{N}{\alpha}\frac{\beta}{\beta-1} \Longleftrightarrow
(m+1)\log\beta \ge W\left(\frac{N}{\alpha}\frac{\beta\log\beta}{\beta-1}\right)\,,
\]
where $W$ is the upper real branch of the Lambert $W$-function, see~\cite{CGHJK96}. We take this branch because $\beta> 1$ so $(m+1)\log\beta > 0$.
Since $W$ is increasing and $\frac{\log \beta}{\beta-1} <1$, this is implied by 
\[
m+1 \ge \frac{W(N\beta/\alpha)}{\log\beta}\,.
\]
Now we have $\log x \ge W(x)$ for all $x\ge e$, so since $\beta>1$ and $N/\alpha \ge e$ this is implied by 
\[
m+1 \ge \frac{\log(N\beta/\alpha)}{\log\beta} = \frac{\log(N/\alpha)}{\log\beta} + 1\,.
\]
Then it suffices to take $m\ge \log(N/\alpha)/\log\beta$. But since $\log\beta\sim\beta-1$ as $\beta\to1$ and $\beta< 1+c$, we can find $c'$ depending only on $c$ such that $c'\log\beta \ge \beta -1$. Then it suffices to take 
\[
m \ge \frac{c'}{\beta-1}\log\left(\frac{N}{\alpha}\right)\,\qedhere
\]
\end{proof}
%
Corollary~\ref{cor:tayloreps} tells us how many terms of a Taylor expansion of $\log g$ we need to get an additive error of at most $\alpha$, under the condition that $g$ has no roots in the disc $\{z\in\CC : |z|<\beta\}$. 
We want to work with a zero-free region of the form $\cN([0,1],\eta)$, the open set containing a ball of radius $\eta$ around each point in $[0,1]$. 
Barvinok~\cite{Barvinok16,Barvinok18} also gives constructions that perform well for this situation which we reproduce below.
%
\begin{lemma}[{\cite[Lemma 2.2.3]{Barvinok16}}]\label{lem:disctostrip}
For $0<\rho<1$ there is a polynomial $p$ of degree
\[
N_\rho := \left\lfloor\left(1+\frac{1}{\rho}\right)e^{1+\frac{1}{\rho}}\right\rfloor \ge 14
\]
such that $p(0)=0$, $p(1)=1$, and 
\[
-\rho \le \Re\, p(z) \le 1 + 2\rho 
\quad\text{and}\quad
|\Im\, p(z)| \le 2\rho
\]
for all $z$ such that $|z|\le \beta(\rho)$ where
\[
\beta_\rho := \frac{1-e^{-1-\frac{1}{\rho}}}{1-e^{-\frac{1}{\rho}}} > 1\,.
\]
\end{lemma}
%
\begin{corollary}[{cf.~\cite[Theorem 1.6]{Barvinok18}}]\label{cor:stripaccuracy}
Suppose that $0<\eta<1$, and $g$ is a polynomial of degree $N$ such that $g(z)\ne0$ for all $z\in\cN([0,1],\eta)$. 
Then given any $0<\alpha<Ne^{6/\eta-1}$, it suffices to compute the first
\[
m=e^{6/\eta}\log\left(\frac{Ne^{6/\eta}}{\alpha}\right)
\]
coefficients of $g$ to obtain a number $\xi$ satisfying $|\log g(1)-\xi| \le \eps$. 
\end{corollary}
\begin{proof}
Let $\rho=\eta/\sqrt 8$ and $\beta_\rho,N_\rho$ be given by Lemma~\ref{lem:disctostrip}. 
Then since $\eta<1$, we note that
\[
N_\rho\le e^{6/\eta}
\quad\text{and}\quad
\beta_\rho \ge 1 + \frac{1}{2e^{\frac{1}{\rho}}} \ge 1+ e^{-4/\eta}\,.
\]
Now the polynomial $p$ as in Lemma~\ref{lem:disctostrip} maps $\{z\in\CC : |z|\le \beta_\rho\}$ into $\cN([0,1], \eta)$, so the polynomial $g\circ p(z)$ is a degree $NN_\rho \le Ne^{6/\eta}$
polynomial which is nonzero for all $z\in\CC$ such that $|z|\le 1+ e^{-4/\eta}$.
%
We now apply Corollary~\ref{cor:tayloreps} to $g\circ p$. 
With $f(z) = \log(g\circ p(z))$ we have $f(1) = \log g(1)$ since $p(1)=1$, and so the Taylor polynomial $T_m(1)$ (as defined in Definition~\ref{def:taylor}) for this $f$ is the quantity we want for $\xi$.
More precisely, as described in~\cite[Section 2.2.2]{Barvinok16}, to compute $T_m(\log g\circ p)$ at $z=0$ it suffices to compute $T_m(g\circ p)$. 
In turn, to compute $T_m(g\circ p)$ it suffices to compute $T_m(g)$ and the truncation $p_m$ of $p$ obtained by deleting all monomials of degree higher than $m$, and then the composition of polynomials $T_m(g) \circ p_m$.
The final step is to truncate $T_m(g) \circ p_m$ to be degree $m$, and to obtain $\xi$ by evaluating this polynomial at $z=1$.
This can be done in time $O(m)$, and by Corollary~\ref{cor:tayloreps} applied to $g\circ p$, when we have
\[
m \ge e^{6/\eta}\log\left(\frac{Ne^{6/\eta}}{\alpha}\right)\,,
\]
and $\alpha < Ne^{6/\eta-1}$, we have the desired $|\log g(1) - \xi| \le\alpha$.
\end{proof}
%
%
We know now how many coefficients are needed for an approximation to a polynomial. 
For the complexity of computing these coefficients we refer to Patel and Regts~\cite{PR17,PR19}.
Our partition function $Z(G;w)$ is an edge-coloured BIGCP in the sense of Patel and Regts' definition in~\cite{PR19}.
$Z(G;w)$ has degree at most $\Delta n/2$ with BIGCP parameters $\alpha=2$ and $\beta_i=O(k^i)$ according to~\cite[Section 6]{PR17}.
Then by~\cite[Theorem 2.1]{PR19} there is a deterministic algorithm to compute the first $m$ coefficients of $Z$ in time
\[
\tilde O\big(n(4e\Delta\sqrt k)^{2m}\big)\,,
\]
where $\tilde O$ means that we omit factors polynomial in $m$.
%
To prove Theorem~\ref{thm:alg,interp} we want an approximation for $\log Z(G;w)$ with error $\alpha=C n$, and we have a zero-free region surrounding $[1,w^*]$ with $w^*=1+(\log k-1)\Delta$ at distance $\eta=\omega(1/(\Delta\log k))$. 
Then we can transform $Z$ into a polynomial with zero-free region around $[0,1]$ of distance $\eta/(w^*-1) = \omega(1/(\log k)^2)$, so we need
\[
C  \le \frac{2}{e\Delta}e^{O(\log^2 k)} \,,
\]
and
\[
m = e^{O(\log^2 k)}\log\left(\frac{\Delta}{2C}e^{O(\log^2 k)}\right)\,,
\]
according to Corollary~\ref{cor:stripaccuracy}.
%
Then we have a running time of
\[
\tilde O\big(n(4e\Delta\sqrt k)^{2m}\big)
= n\exp\left(e^{O(\log^2 k)} \log\left(\frac{\Delta}{C}e^{O(\log^2 k)}\right)\log(\Delta\sqrt k)\right)\,.
\]
%
